\documentclass[11pt]{article}
\usepackage[margin=1in]{geometry}
\usepackage{amsmath, amssymb}
\usepackage{booktabs}
\usepackage{hyperref}

\title{World-Model KL Surprise vs.\ C51 Mutual Information\\
\large Two intrinsic signals for the Light-Dark POMDP}
\author{}
\date{}

\begin{document}
\maketitle

\paragraph{Sign convention.}
WM-KL is \emph{high} inside the band, not low. In the band, the precise
observation makes the posterior $q(z_{t+1}\mid x_{t+1}, z_t, a_t)$ sharpen
far from the broader prior $p(z_{t+1}\mid z_t, a_t)$, so the KL is large. In
the dark, the observation is too noisy to update the prior much, so the KL is
small. Both WM-KL and C51-MI therefore point the agent toward the band; they
just do so through different mechanisms.

\paragraph{One-line framing.}
\begin{itemize}
  \item \textbf{WM-KL:} ``I learned something about the world.''
  \item \textbf{C51-MI:} ``I learned something that matters for the task.''
\end{itemize}

\section{What each signal measures}

\subsection*{WM-KL --- state-representation level}
\[
  r_{\mathrm{WM\text{-}KL}}(t{+}1)
  \;=\;
  \mathrm{KL}\!\left(
    q(z_{t+1}\mid x_{t+1}, z_t, a_t)
    \;\big\|\;
    p(z_{t+1}\mid z_t, a_t)
  \right).
\]
This is the information a single observation provided about where the agent
is, regardless of what that location means for the task. It is:
\begin{itemize}
  \item per-step (a snapshot, not a difference);
  \item task-agnostic --- the reward function never enters the formula;
  \item driven by \emph{observation noise structure} (low $\sigma$ $\Rightarrow$ large KL; high $\sigma$ $\Rightarrow$ small KL);
  \item about the \emph{world model's} own uncertainty over the next latent.
\end{itemize}

\subsection*{C51-MI --- return-distribution level}
\[
  r_{\mathrm{info}}(t)
  \;=\;
  \mathrm{MI}(t) - \mathrm{MI}(t{+}1),
  \qquad
  \mathrm{MI}(t) \;=\; I\!\left(Z;\, R \mid a\right)
  \;=\;
  H\!\left[\bar{p}(R\mid a)\right]
  - \mathbb{E}_z\!\left[ H\!\left[p(R\mid z, a)\right] \right].
\]
This measures how much, in expectation, knowing the latent reduces
uncertainty about future return under the current belief. It is:
\begin{itemize}
  \item a temporal difference (sharpening over a step);
  \item task-conditional --- the return distribution is shaped by reward and~$\gamma$;
  \item driven by how disambiguating different latents are for outcomes;
  \item about the agent's \emph{predictive uncertainty over consequences}.
\end{itemize}

\section{The conceptual axis}

Consider four kinds of transitions:

\begin{center}
\begin{tabular}{clccc}
\toprule
 & Description & Observation informative? & Matters for return? \\
\midrule
A & Band entry in Light-Dark        & yes & yes \\
B & Informative but task-irrelevant feature & yes & no  \\
C & State-level redundant but task-relevant disambiguation & no  & yes \\
D & Empty rollout, deep in the dark & no  & no  \\
\bottomrule
\end{tabular}
\end{center}

\begin{center}
\begin{tabular}{ccc}
\toprule
Case & WM-KL & $r_{\mathrm{info}}$ \\
\midrule
A & high & high \\
B & high & $\approx 0$ \\
C & $\approx 0$ & high (if belief was wrong) \\
D & $\approx 0$ & $\approx 0$ \\
\bottomrule
\end{tabular}
\end{center}

In Light-Dark, cases A and D dominate, so the two signals are largely
aligned. The philosophical difference, however, persists:

\begin{itemize}
  \item \textbf{WM-KL} says: \emph{``seek observation-rich states.''} The
    agent learns the structure of the world, full stop.
  \item \textbf{C51-MI} says: \emph{``seek states whose disambiguation
    tightens your prediction of return.''} The agent learns the structure
    of \emph{the task as it applies to its current belief}.
\end{itemize}

\section{Why C51-MI is, in principle, the more refined signal}

C51-MI carries information that WM-KL discards:

\begin{enumerate}
  \item \textbf{Task gating.} WM-KL fires wherever observations are
    informative. In a multi-room maze with cameras on every wall, WM-KL
    would happily reward staring at wall textures. C51-MI fires only when
    disambiguation actually changes the return prediction.
  \item \textbf{Reducibility (noisy-TV problem).} WM-KL does not know which
    uncertainty is reducible by agent action. A noisy but uncontrollable
    observation channel will keep producing high WM-KL forever. C51-MI is
    grounded in the return distribution, which only sharpens when the
    agent's value estimate tightens---a much more action-controllable
    quantity.
  \item \textbf{Belief-conditional priority.} WM-KL is independent of what
    the agent currently believes. C51-MI uses the actual belief at time~$t$.
    If the agent already knows where it is in the band, both $\mathrm{MI}(t)$
    and $\mathrm{MI}(t{+}1)$ are small and $r_{\mathrm{info}} \approx 0$;
    redundant band visits are not paid for. WM-KL would still pay.
\end{enumerate}

\noindent
In short: \textbf{C51-MI is WM-KL filtered through ``and this matters for
what I'm trying to do.''}

\section{Why WM-KL is, in practice, the more usable signal}

The flip side is bootstrapping. The refinement of C51-MI comes from passing
through the return distribution $p(R\mid z, a)$, which is only meaningful
once the C51 head is trained. Early in training:
\begin{itemize}
  \item The C51 head is near-uniform $\Rightarrow$ $I(Z;R)$ is whatever
    noise the head outputs $\Rightarrow$ $r_{\mathrm{info}}$ is noise.
  \item WM-KL only requires the world-model encoder/decoder/prior to be
    plausible, which holds after a few thousand world-model gradient steps,
    independent of the critic.
\end{itemize}
Thus WM-KL is \textbf{available earlier}, while C51-MI is \textbf{better-shaped later}.

\section{Complementarity}

The two signals form a natural two-phase decomposition of exploration:

\begin{itemize}
  \item \textbf{WM-KL is the natural curriculum for the encoder.} Driving
    the agent to observation-rich states gives the world model the data
    it needs to become accurate. It is a bootstrapping signal: it improves
    the substrate every other component---including C51---relies on.
  \item \textbf{C51-MI is the natural objective for the policy.} Once the
    world model is good enough that beliefs are meaningful, C51-MI says
    ``go to states whose disambiguation will tighten your value estimate.''
    It directly aligns belief reduction with task progress.
\end{itemize}

Sequenced informally:
\[
  \text{Phase 1: WM-KL dominates} \;\longrightarrow\; \text{agent learns the world's observation structure.}
\]
\[
  \text{Phase 2: C51-MI dominates} \;\longrightarrow\; \text{agent learns which parts of that structure matter.}
\]

This mirrors the strategy of methods like Plan2Explore (disagreement-based,
encoder-driven) followed by task-aware refinement once reward starts flowing.

\section{Concrete proposal}

Combine the two as a weighted shaping term:
\[
  r_{\mathrm{total}}(t)
  \;=\;
  r_{\mathrm{env}}(t)
  \;+\;
  \beta_{\mathrm{KL}}(t)\cdot r_{\mathrm{WM\text{-}KL}}(t)
  \;+\;
  \beta_{\mathrm{MI}}(t)\cdot r_{\mathrm{info}}(t),
\]
with $\beta_{\mathrm{KL}}(t)$ annealed downward and $\beta_{\mathrm{MI}}(t)$
annealed upward over training. The existing z-score normalization and
adaptive-$\beta$ controller can be reused unchanged for both terms.

\paragraph{Recommended first experiment.}
Try WM-KL alone first. The current C51-MI runs show the C51 head's
belief-entropy proxy is not yet a reliable signal source---it actively hurts
$\mathrm{success\_reach}$ at active magnitudes. WM-KL would test whether
exploration shaping helps \emph{at all}, using a signal that does not
depend on the critic being calibrated.
\begin{itemize}
  \item If WM-KL alone moves $\mathrm{success\_reach}$ meaningfully, the
    bottleneck is exploration; C51-MI can be layered back on later.
  \item If WM-KL also does not move it, the bottleneck is value
    propagation rather than exploration, and the right next steps are
    $n$-step returns, a separate intrinsic value head, or simply longer
    training.
\end{itemize}

\end{document}
